\begin{lemma}[The original projective integral global function ring]
\label{lem:fc-original-projective-integral-global-functions}
\lean{EnvelopeProperty.ProjectiveIntegralGlobalFunctions.sections_finite, EnvelopeProperty.ProjectiveIntegralGlobalFunctions.sections_isField}
\leanok
For the original integral projective $p:X\to\operatorname{Spec}k$, the
entire actual ring $\Gamma(X,\mathcal O_X)$, with its original
structure-map coefficient action, is a finite field extension of $k$.
\end{lemma}
\begin{proof}
The original structure line is finite presentation and quasicoherent.
Actual projective coherent finiteness gives finite-dimensional original
field-linear global sections.  The actual whole structure-section
isomorphism identifies these with $\Gamma(X,\mathcal O_X)$ preserving
every original coefficient.  Original integrality makes this nonzero
ring a domain.  A domain finite dimensional over a field is a field.
\end{proof}

\begin{lemma}[Nonzero original global functions are whole units]
\label{lem:fc-original-projective-global-function-units}
\lean{EnvelopeProperty.ProjectiveIntegralGlobalFunctions.section_isUnit, EnvelopeProperty.ProjectiveIntegralGlobalFunctions.unitEndomorphism_isIso}
\leanok
\uses{lem:fc-original-projective-integral-global-functions}
Every nonzero original global function is an actual global unit.  Its
original multiplication morphism is an automorphism of the whole
original structure line.
\end{lemma}
\begin{proof}
Use the preceding actual finite-dimensional domain to obtain an inverse
in the same original global ring.  Multiplication by this original unit
and its inverse gives the canonical whole structure-module automorphism.
\end{proof}
